\section{怎样辨析词类}

大家知道，现代汉语的词，按照它们的不同意义和在造句时所起的作用不同，可以分为实词和虚词两大类。实词具有实在的意义，能够作句子成分，并且在一定的话语环境里能够独立成句，回答一定问题；虚词则不表示实在意义，它们的基本作用是表示语法关系，一般不能单独作句子成分。
\newpage
\noindent
（副词例外），不能单独回答问题（极少数副词“不”、“没有”、“也许”等例外）。实词包括名词、动词、形容词、数词、量词、代词等六类；虚词包括副词、介词、连词、助词、叹词等五类。所以说，现代汉语词类分实词和虚词两个大类和一个小类。

那么，我们怎样才能确定句子中的某个词是属于哪类词呢？

1. 辨析一个词能否表示事物及事物的动作、变化、性质、状态、数量等概念，能否单独回答问题。能单独回答问题的，是实词；不能单独回答问题的，是虚词（少数副词能单独回答问题，但是不能表示事物及其有关的概念，仍属虚词）。例如：

甲：喂，你这两本新书借给我看看，好吗？

乙：可以。不过我刚从图书馆借来，还没看，请你快一点还回来。

这两句对话中，除了加点的字没有实在的意义，也不能单独回答问题（副词“没”能单独回答问题，但无实义）之外，别的词都有比较实在的意义，并能单独回答问题。据此可以划分虚词和实词。“喂”是一个叹词，常独立于句子之外，是一种特殊的虚词；其余的“刚、还hái”（副词），“从”（介词），“不过”（连词），“吗”（语气助词）都只能附着在词、词组、句子上面，起帮助造句的语法作用，所以它们都是虚词。而“借、给、看、请、还huán、回、来、可以”，表示事物的动作变化的概念（动词），“书、图书馆”等词表示事物的名称（名词）；“新、好、快”等词表示事物的性质状态的概念（形容词）；“两本、一点”等
\newpage
\noindent
词表示事物数量的概念（数词、量词），“你、我、这”等词表示代替有关实词的概念（代名词），而且都可以单独回答问题，所以它们都是实词。

2.掌握各类实词和各类虚词的代表词，可以由此类推相应的各类实词。汉语的词语极为丰富，要逐一地记住它们的词类是不可能的，也没有必要。比如名词吧，就有表示人的名称的“工人、妇女、鲁迅、解放军、人类、侏儒”等等；表示事物名称的“房屋、船只、熊猫、石头、玫瑰、拖拉机”等等；表示抽象事物名称的“政治、科学、智慧、思维、概念、创造性”等等；表示地点处所名称的“北京、边疆、中国、亚洲、解放路、天安门”等等；表示时间名称的“今天、早晨、中午、晚上、现在、星期日”等等。再如副词吧，就有表示程度的“很、极、非常、格外、越、稍微”等等；表示范围的“都、全、统统、只、仅仅、一概”等等；表示时间的“已、曾、刚、才、正、将、立刻”等等；表示频率的“再、常、一向、终于、突然、渐渐、仍旧”等等；表示肯定、否定、可能的“必、准、不、没、未必、或许、大约”等等；表示语气的“可、却、倒、竟、偏、岂、难道、其实、未免”等等。这真是数不胜数，记不胜记。我们可以在了解一个词类所具有的几方面的意义的基础上，选择能代表某方面意义的代表词，以此类推，就可以大致解决一个词究竟属于什么词类的问题了。我们可以把这些代表词列表如下，请加以记忆和类推。（见下页）

比如我们记到了动词的几个代表词，就可由“唱”想到一系列表示动作的词，由“宣传”想到一系列表示行为的词，由“发展”想到一系列表示变化的词，由“希望”想到
\newpage
\noindent
\begin{tabular}{p{2em}|p{2em}|p{1em}|p{4em}|p{5cm}}
	\hline
	\multicolumn{4}{c|}{词类} & 代\quad表\quad词\\
	\hline
	实词 & 名词 & \multicolumn{2}{c|}{名词} & 工人、房屋、政治、北京、今天\\
	\hline
	实词 & 名词 & 附: & 方位词 & 前、上、东面、中间\\
	\hline
	实词 & 动词 & \multicolumn{2}{|c|}{动词} & 唱、宣传、发展、希望、存在、有\\
	\hline
	实词 & 动词 & 附: & 能愿动词 & 能、敢、应该、可以、必须、愿意\\
	\hline
	实词 & 动词 & 附: & 趋向动词 & 来、上、开、进、过、下去、回来\\
	\hline
	实词 & 动词 & 附: &判断词 & 是\\
	\hline
	实词 & \multicolumn{3}{|c|}{形容词} & 好、勇敢、平常；大、红、畅快、干净\\
	\hline
	实词 & \multicolumn{3}{|c|}{数词} & 一、十、百、千、万\\
	\hline
	实词 & \multicolumn{3}{|c|}{量词} & 个、件、次、趟\\
	\hline
	实词 & 代词 & \multicolumn{2}{|c|}{人称代词} & 我、它、她们、自己、大家\\
	\hline
	实词 & 代词 & \multicolumn{2}{|c|}{疑问代词} & 谁、什么、哪、怎样、多少\\
	\hline
	实词 & 代词 & \multicolumn{2}{|c|}{指示代词} & 这、那儿、这会儿、那些、每、某\\
	\hline
	虚词 & \multicolumn{3}{|c|}{副词} & 很、都、已经、的确、不、竟、却、再\\
    \hline
    虚词 & \multicolumn{3}{|c|}{介词} & 向、在、对、跟、当、把、被、比、因\\
    \hline
    虚词 & \multicolumn{3}{|c|}{连词} & 和、或者、以及、但是、所以\\
    \hline
	虚词 & 助词 & \multicolumn{2}{|c|}{结构助词} & 的、地、得、所\\
	\hline
	虚词 & 助词 & \multicolumn{2}{|c|}{时态助词} & 着、了、过\\
    \hline
	虚词 & 助词 & \multicolumn{2}{|c|}{语气助词} & 呢、吗、吧、哇(wa)\\
    \hline
    虚词 & \multicolumn{3}{|c|}{叹词} & 啊、唉、扔、嗯、哈哈\\
    \hline
\end{tabular}	
\newpage

一系列表示心理活动的词，由“存在”想到表示存在的词，由“有”想到表示有无的词，而这些词都是动词。其他名词类的代表词，也各自代表某一方面的意义或语法作用，我们只要了解某一词类具有的几方面意义或作用，就可以由这些代表字类推出同性质的一类词来。

3. 根据各类词的基本语法特点，来辨析一个词属于哪类词。例如，名词的一个明显的语法特点，就是受数量词的修饰，不受副词的修饰（只有极少数抽象名词，如“科学、道德、标准”等可受副词修饰），我们可以说“一位老师，五本书，八间教室”，而不能说“很老师，不书，已经教室”；动词则可受时间副词的修饰，而除了表示心理活动的动词外，其他单个的动词一般都不受程度副词的修饰，我们可以说“正在进行，立刻讨论，已经休息”，而不能说“很进行，非常讨论，极休息”；形容词在这一点上与动词相反，一般能受程度副词的修饰，而一般不受时间副词的修饰，我们可以说“很高大，极遥远，非常丰富”，而不说“正在高大，立刻遥远，已经丰富”，（当单个的形容词带上别的词后，一般就可以受时间副词的修饰了，如“正在高大起来，立刻很远远了，已经很丰富了”。）

另外，一般的动词都能带上相应的宾语，而一般的形容词不能带宾语。例如：我们常说“大家讲道理”，而不说“大家大道理”（“大”是形容词，只能修饰“道理”），我们常说“我们学习语法”，而不说“我们正确语法”（“正确”是形容词，只能修饰“语法”）。动词后边可以带时态助词和趋向动词表示变化；而介词则不能，只可以与名词、代词或词组等构成介词结构作句子成分，自身不能独立作为句子成分。例如：
\newpage

“我们按照规定执行了。”

在这儿，介词“按照”后面不能附加“了、着、过”等时态助词或“来、去、上来、下去”等趋向动词，只能与名词“规定”构成介词结构去修饰动词“执行”；而动词“执行”后面可以加时态助词或趋向动词，充当句子的谓语成分。

当然，名词、形容词、动词、介词等都还有其他的语法特点，别类词也同样有自己的语法特点（如副词不能修饰名词，只能修饰动词或形容词；助词只能附着在词、词组或句子上表示结构关系，时态变化或者种语气，等等），我们初学词类，没有必要将它们的一切特点完全辨析清楚，但是，这些词类的基本语法特点却应该掌握，因为了解了这些基本语法特点，不仅能很好地帮助我们分清词类，而且有利于辨析和避免由于不明词性而造成的病句。例如下面一段话中加点的词，究竟属于哪类词呢？

“我们（ ）生活在一个（ ）开辟人类新历史的光辉时代（ ）。在（ ）这样的（ ）时代，人们对（ ）许许多多的事物都产生了新的联想、新的感情。不是（ ）有许多人在（ ）讴歌那（ ）光芒四射的朝阳、四季常青的松柏、庄严屹立（ ）的山峰、澎湃（ ）翻腾的海洋吗？不是有好些（ ）人在赞美挺拔（ ）的白杨、明亮的灯火、奔驰（ ）的列车、崭新的日历吗（ ）？睹物思人，这些东西（ ）引起人们多少（ ）丰富和（ ）充满感情的想象（ ）！”

\hspace{13em}
-----秦牧《土地》
\newpage

需要注意的是：

① “在这样的时代”的“在”与“不是有许多人在讴歌……吗”的“在”（“正在”的意思），只是同音同形而意义毫无联系的同音词而已（“父母健在”中的“在”是一个表示“存在、生存”意思的动词，则是另一个同音词），不可混为一谈。前者是介词，后者是副词。

② “新的联想”中的“联想”，“充满感情的想象”中的“想象”，都具有名词和动词两种词性，词性虽然不同，但词汇意义没有截然的不同，从词类角度讲，就叫作兼类。在这儿，它们都是名词。

我们在辨清词类的基础上，就可以避免或修改由于不明词类而造成的病句。例如：

(1) 我班夺得了流动红旗，同学们都感到非常高兴和荣誉。
(“荣誉”是名词，不受副词“非常”修饰，也不能与形容词“高兴”并列。应改为形容词“光荣”。)

(2) 尽管有人还在流言语，然而他那改革的决心毫不动摇。
(“流言语”是个成语，指没有根据的话，属于名词性词组，不受副词“还在”的修饰。这儿误把名词性词组当作动词使用了。应该在“流言语”之前，加上一个动词“散布”或“制造”。)

(3) 小说《高山下的花环》很感动人，能干净我们的思想。
(“干净”是形容词，而一般形容词都不能带宾语，这儿将形容词误用为动词了。应将“干净”改为动词“净化”。）

\newpage
（4）现在已经多数人来了，就开会吧。
（“已经”是时间副词，而副词不能修饰名词，只能修饰动词或形容词。因此，“已经”要移到动词“来”之前。）

（5）他不求上进，你不该跟着他比。
（“跟”兼属动词、介词两类。在这儿是介词，与“他”这个代词构成介词结构，充当动词“比”的状语。而介词是不能同时态助词“着、了、过”组合的，应该删去“着”。）

我们只要了解了关于现代汉语词类的基本知识和辨析词类的一些方法，进行必要的练习，那么，辨析汉语词类的问题也是不难解决的。